% ---- Abstract (EN) + Zusammenfassung (DE) --------------------------------------------
\chapter*{Abstract}
\addcontentsline{toc}{chapter}{Abstract}

Predictive Coding (PC) is a biologically motivated, local-learning alternative to
backpropagation (BP) whose reported advantages are inconsistent across the literature. This
thesis builds a from-scratch State-Optimization PC network on MNIST in pure PyTorch---no
autograd in the learning rule; the hand-derived settling and weight gradients are verified
against autograd---and uses it for a deliberately sober, reproducible study with three
contributions. \textbf{(i)~Systems:} among open PC libraries, none ships a custom inference
kernel; we contribute, to our knowledge, the first hand-written \emph{fused CUDA settling
kernel} for PC, show that PC inference at MNIST scale is launch-overhead-bound, and obtain
$3.2\times$ at batch~64, generalizing the kernel to arbitrary depth and using it as the engine
of one of our own studies ($1.45\times$ end-to-end, identical results). \textbf{(ii)~Methodology:}
a confound-controlled fair PC-vs-BP comparison (matched architecture, initialization, data and
loss; per-method learning-rate tuning; bootstrap confidence intervals over multiple seeds) finds
\emph{PC\,$\approx$\,BP} across accuracy, robustness, sample efficiency and continual learning,
including an architecture-faithful replication of Song et~al.'s exact alternating
Split-FashionMNIST protocol where no PC-vs-BP difference exceeds one standard deviation at any
training budget. We isolate three confounds (loss function, the stability--plasticity trade-off,
and an untuned baseline) that manufacture spurious ``PC advantages'', and describe an adversarial
verification protocol that caught a literature-framing error before it entered this document.
\textbf{(iii)~Capability:} a generatively trained PC recovers the ``one mechanism, many tasks''
property---classification, generation, inpainting, and anomaly detection from a single settling
rule. We additionally implement incremental PC and an EWC baseline. All code, tests, figures, and
experiment artifacts are released publicly.

\vspace{1cm}
\begin{otherlanguage}{ngerman}
\section*{Zusammenfassung}

Predictive Coding (PC) ist eine biologisch motivierte, lokal lernende Alternative zur
Backpropagation (BP), deren berichtete Vorteile in der Literatur uneinheitlich sind. Diese
Arbeit implementiert ein From-Scratch-PC-Netz (State-Optimization-Formulierung) auf MNIST in
reinem PyTorch---ohne Autograd in der Lernregel; die handabgeleiteten Gradienten sind gegen
Autograd verifiziert---und nutzt es f\"ur eine bewusst n\"uchterne, reproduzierbare Studie mit
drei Beitr\"agen. \textbf{(i)~Systems:} der erste handgeschriebene, fusionierte CUDA-Settling-Kernel
f\"ur PC (nach unserem Kenntnisstand); PC-Inferenz ist auf MNIST-Skala launch-overhead-gebunden,
der Kernel erreicht $3{,}2\times$ bei Batch~64, ist auf beliebige Tiefe verallgemeinert und treibt
eine der Studien ($1{,}45\times$, identische Ergebnisse). \textbf{(ii)~Methodik:} ein
confound-kontrollierter, fairer PC-vs-BP-Vergleich findet \emph{PC\,$\approx$\,BP}, inklusive einer
architektur-treuen Replikation von Songs exaktem alternierenden Protokoll (kein Unterschied
$>1\sigma$ bei keinem Budget). Drei Confounds (Loss-Funktion, Stabilit\"at-Plastizit\"at,
ungetunte Baseline) erzeugen scheinbare PC-Vorteile; eine adversariale Verifikation fing einen
Framing-Fehler vor Eingang in diese Arbeit. \textbf{(iii)~F\"ahigkeit:} ein generativ trainiertes
PC liefert die ``eine Maschinerie, viele Aufgaben''-Eigenschaft. Zus\"atzlich sind inkrementelles
PC und eine EWC-Baseline implementiert. Aller Code, Tests, Figuren und Artefakte sind
ver\"offentlicht.
\end{otherlanguage}
